% intro.tex — Expanded ICLR introduction.

\section{Introduction}
\label{sec:intro}

Modern protein machine learning typically represents proteins as amino-acid
sequences or atomic coordinates. Large protein language models such as ESM and
ESM-2 learn evolutionary and functional patterns directly from sequence
\citep{rives2021biological,lin2023evolutionary}. Structure predictors such as
AlphaFold and ESMFold infer accurate atomic models from sequence
\citep{jumper2021highly,lin2023evolutionary,abramson2024accurate}. Together,
these representations support protein function prediction, binding-affinity
estimation, and structure-based drug design
\citep{radivojac2013large,ozturk2018deepdta,pinheiro2024structure}.

Coordinates, however, do not preserve all experimental evidence behind a
crystallographic structure. An X-ray structure is a model refined against a
diffraction-derived electron-density map \citep{adams2011phenix}. A single
coordinate set may omit weak evidence for alternative conformations
\citep{lang2010automated,lang2014protein}. The map also indicates how strongly
individual atoms and molecular fragments are supported by the data
\citep{meyder2017estimating}. Coordinate-only inputs therefore lose direct
access to this experimental evidence.

This evidence may be especially relevant at protein--ligand interfaces, where
small structural changes alter intermolecular contacts. Electron-density
analysis has revealed unmodeled ligand conformations and alternative binding
modes \citep{van2018qfit}. Prior work has also used experimental density for
ligand identification, affinity prediction, and molecular generation
\citep{kowiel2019automatic,suriana2025beyond,wang2022pocket}. These approaches
integrate density into task-specific pipelines. Whether a shared,
density-pre-trained representation can transfer across distinct
protein--ligand tasks remains unclear.

We introduce \ours{}, a coordinate-and-density encoder for protein--ligand
complexes. As illustrated in \cref{fig:method-overview}, it places ligand atoms,
pocket atoms, experimental density, and its gradient magnitude on a common
voxel grid. We curate more than 100K crystallographic protein--ligand instances
from PLINDER \citep{durairaj2024plinder}. We then pre-train a Channel-ViT through
masked reconstruction \citep{bao2023channel-24a,he2022masked}. At transfer time,
the frozen encoder provides either pooled complex embeddings or spatial
features.

We test whether these representations support both prediction and generation.
Binding-affinity regression examines whether pooled features capture
interaction strength. Reference-assisted molecular generation examines whether
spatial features improve ligand construction within a target pocket. We
evaluate affinity under complex-level leakage and protein-sequence-similarity
constraints, including an exact-PDB-held-out CASF-2016 subset. For generation,
we examine docking, molecular quality, reference-ligand similarity, pose
geometry, and qualitative binding modes.

To summarize, our contributions are as follows:
\begin{itemize}
    \item We introduce \ours{} and a large-scale curation pipeline for joint
    pre-training on fitted coordinates and experimental electron-density maps.
    \item We demonstrate the utility of the frozen representation for
    binding-affinity regression under complex-level leakage and
    protein-sequence-similarity controls.
    \item We transfer the same encoder to molecular generation and evaluate
    docking, molecular quality, reference similarity, and pose geometry.
\end{itemize}
